\documentclass[conference, a4paper]{IEEEtran}

% Character set and language settings for English
\usepackage[utf8]{inputenc}
\usepackage[T1]{fontenc}
\usepackage[english]{babel}

% Useful packages for IEEE papers
\usepackage{cite}
\usepackage{amsmath,amssymb,amsfonts}
\usepackage{algorithmic}
\usepackage{graphicx}
\usepackage{textcomp}
\usepackage{xcolor}

\begin{document}

\title{Title of the IEEE Paper in English\\
{\footnotesize \textsuperscript{*}Note: Sub-titles or project names can go here}}

\author{
\IEEEauthorblockN{1. First Name Last Name}
\IEEEauthorblockA{\textit{Department / Faculty} \\
\textit{Name of University or Company}\\
City, Country \\
Email address}
\and
\IEEEauthorblockN{2. First Name Last Name}
\IEEEauthorblockA{\textit{Department / Faculty} \\
\textit{Name of University or Company}\\
City, Country \\
Email address}
}

\maketitle

\begin{abstract}
This document is a template for an IEEE conference paper in English. The abstract should briefly summarize the problem statement, the methodology applied, as well as the main results and conclusions of the paper.
\end{abstract}

\begin{IEEEkeywords}
IEEE format, LaTeX, template, English, seminar
\end{IEEEkeywords}

\section{Introduction}
This document demonstrates how to use the classic IEEE layout for papers in English. The format defines a two-column layout, specific heading styles, and guidelines for references \cite{b1}.

\section{Methodology}
This section typically describes the approach or system model. You can insert mathematical formulas here, such as Equation \eqref{eq:example}.

\begin{equation}
E = mc^2 \label{eq:example}
\end{equation}

Bullet points can also be used:
\begin{itemize}
    \item First important point.
    \item Second important point regarding the methodology.
\end{itemize}

\section{Results}
Present your results here. Charts and tables should be centered. Table captions are always placed above the table in the IEEE format.

\begin{table}[htbp]
\caption{Example Table}
\begin{center}
\begin{tabular}{|c|c|c|}
\hline
\textbf{Column 1} & \textbf{Column 2} & \textbf{Column 3} \\
\hline
Value A & Value B & Value C \\
\hline
\end{tabular}
\label{tab1}
\end{center}
\end{table}

\section{Conclusion}
The conclusion summarizes the findings of the paper and may provide an outlook on future research opportunities or extensions.

\begin{thebibliography}{00}
\bibitem{b1} G. Eason, B. Noble, and I. N. Sneddon, ``On certain integrals of Lipschitz-Hankel type involving products of Bessel functions,'' Phil. Trans. Roy. Soc. London, vol. A247, pp. 529--551, April 1955.
\end{thebibliography}

\end{document}